% ============================================================
% Question Bank: ch09-03 Experimental Probability
% Topic Title: Experimental Probability
% CCSS: 7.SP.C.6
% Questions: 27 (15 MC + 6 SA + 6 Graphical)
% ============================================================

% --- Multiple Choice Questions (15) ---

\begin{practiceQuestion}{9-3-q01}{mc}
\begin{questionText}
A factory inspects $250$ tablets and finds that $15$ have screen defects. What is the experimental probability that a randomly selected tablet has a screen defect?
\end{questionText}
\choiceA{$\dfrac{3}{50}$}
\choiceB{$\dfrac{15}{250}$}
\choiceC{$\dfrac{1}{15}$}
\choiceD{$\dfrac{3}{25}$}
\correctAnswer{A}
\explanation{Experimental probability equals observed occurrences divided by total trials. With $15$ defects in $250$ tablets: $\frac{15}{250} = \frac{3}{50}$. Choices A and B both equal $\frac{3}{50}$, but A is in simplified form.}
\end{practiceQuestion}

\begin{practiceQuestion}{9-3-q02}{mc}
\begin{questionText}
Amir rolls a standard number cube $300$ times. He records that the number $4$ appears $58$ times. What is the experimental probability of rolling a $4$?
\end{questionText}
\choiceA{$\dfrac{1}{6}$}
\choiceB{$\dfrac{58}{300}$}
\choiceC{$\dfrac{29}{150}$}
\choiceD{$\dfrac{4}{300}$}
\correctAnswer{C}
\explanation{Experimental probability is $\frac{\text{times event occurred}}{\text{total trials}} = \frac{58}{300} = \frac{29}{150}$.}
\end{practiceQuestion}

\begin{practiceQuestion}{9-3-q03}{mc}
\begin{questionText}
A bakery tracks that $36$ out of $180$ customers order chocolate cake. Based on this data, what is the experimental probability that the next customer orders chocolate cake?
\end{questionText}
\choiceA{$\dfrac{1}{4}$}
\choiceB{$\dfrac{1}{6}$}
\choiceC{$\dfrac{1}{5}$}
\choiceD{$\dfrac{2}{9}$}
\correctAnswer{C}
\explanation{Experimental probability equals $\frac{36}{180} = \frac{1}{5}$. This is the best estimate based on the observed data.}
\end{practiceQuestion}

\begin{practiceQuestion}{9-3-q04}{mc}
\begin{questionText}
A basketball team won $42$ of its last $60$ games. Based on this data, which is the best prediction for the probability that the team wins its next game?
\end{questionText}
\choiceA{$\dfrac{7}{10}$}
\choiceB{$\dfrac{1}{2}$}
\choiceC{$\dfrac{3}{5}$}
\choiceD{$\dfrac{42}{60}$}
\correctAnswer{A}
\explanation{Experimental probability is $\frac{42}{60} = \frac{7}{10}$. While choice D shows the same value unsimplified, the simplified fraction $\frac{7}{10}$ is the standard answer.}
\end{practiceQuestion}

\begin{practiceQuestion}{9-3-q05}{mc}
\begin{questionText}
Emma tracks the color of cars passing her house. Out of $400$ cars, $96$ are white. The theoretical probability in her city is $\dfrac{1}{4}$. How does her experimental probability compare?
\end{questionText}
\choiceA{The experimental probability is equal to the theoretical probability.}
\choiceB{The experimental probability is greater than the theoretical probability.}
\choiceC{The experimental probability is less than the theoretical probability.}
\choiceD{There is not enough information to compare.}
\correctAnswer{C}
\explanation{The experimental probability is $\frac{96}{400} = \frac{24}{100} = 0.24$. The theoretical probability is $\frac{1}{4} = 0.25$. Since $0.24 < 0.25$, the experimental is less than the theoretical.}
\end{practiceQuestion}

\begin{practiceQuestion}{9-3-q06}{mc}
\begin{questionText}
A bottling plant tests $1{,}000$ bottles for leaks and finds $8$ are defective. What is the best estimate for the probability that any given bottle is defective?
\end{questionText}
\choiceA{$\dfrac{1}{1000}$}
\choiceB{$\dfrac{1}{125}$}
\choiceC{$\dfrac{8}{100}$}
\choiceD{$\dfrac{1}{8}$}
\correctAnswer{B}
\explanation{Experimental probability is $\frac{8}{1000} = \frac{1}{125}$. This is the best estimate of the defect rate based on testing.}
\end{practiceQuestion}

\begin{practiceQuestion}{9-3-q07}{mc}
\begin{questionText}
If a student conducts an experiment $10$ times and then repeats it $1{,}000$ times, what is likely true about the experimental probabilities?
\end{questionText}
\choiceA{Both experiments will give the exact same probability.}
\choiceB{The $10$-trial experiment will be more reliable.}
\choiceC{The $1{,}000$-trial experiment will likely give a probability closer to the theoretical probability.}
\choiceD{The number of trials does not affect reliability.}
\correctAnswer{C}
\explanation{The Law of Large Numbers states that as the number of trials increases, experimental probability tends to approach theoretical probability. More trials produce more reliable estimates.}
\end{practiceQuestion}

\begin{practiceQuestion}{9-3-q08}{mc}
\begin{questionText}
A seed company tests $500$ sunflower seeds and finds that $435$ germinate. What is the experimental probability that a sunflower seed will germinate?
\end{questionText}
\choiceA{$\dfrac{87}{100}$}
\choiceB{$\dfrac{43}{50}$}
\choiceC{$\dfrac{9}{10}$}
\choiceD{$\dfrac{7}{8}$}
\correctAnswer{A}
\explanation{Experimental probability is $\frac{435}{500} = \frac{87}{100}$. This means about $87\%$ of seeds are expected to germinate.}
\end{practiceQuestion}

\begin{practiceQuestion}{9-3-q09}{mc}
\begin{questionText}
Alex records the number of text messages he receives each hour for $50$ hours. He receives at least one message in $35$ of those hours. Based on this data, what is the experimental probability he receives at least one message in a random hour?
\end{questionText}
\choiceA{$\dfrac{3}{5}$}
\choiceB{$\dfrac{7}{10}$}
\choiceC{$\dfrac{35}{50}$}
\choiceD{$\dfrac{5}{7}$}
\correctAnswer{B}
\explanation{Experimental probability is $\frac{35}{50} = \frac{7}{10}$. Choices B and C represent the same value, but B is simplified.}
\end{practiceQuestion}

\begin{practiceQuestion}{9-3-q10}{mc}
\begin{questionText}
Mia spins a spinner $120$ times and records the results in the table below.

\begin{center}
\begin{tabular}{|c|c|c|c|}
\hline
Red & Blue & Yellow & Green \\
\hline
$28$ & $34$ & $30$ & $28$ \\
\hline
\end{tabular}
\end{center}

Which color has the highest experimental probability?
\end{questionText}
\choiceA{Red}
\choiceB{Blue}
\choiceC{Yellow}
\choiceD{Green}
\correctAnswer{B}
\explanation{Blue appeared the most times ($34$ out of $120$). Since experimental probability is based on observed frequency, blue has the highest experimental probability of $\frac{34}{120} = \frac{17}{60}$.}
\end{practiceQuestion}

\begin{practiceQuestion}{9-3-q11}{mc}
\begin{questionText}
A restaurant tracks that $72$ out of $240$ meals ordered are pasta dishes. Based on this data, how many pasta dishes should the restaurant expect to sell in the next $400$ orders?
\end{questionText}
\choiceA{$72$}
\choiceB{$120$}
\choiceC{$100$}
\choiceD{$160$}
\correctAnswer{B}
\explanation{The experimental probability is $\frac{72}{240} = \frac{3}{10}$. For $400$ orders: $\frac{3}{10} \times 400 = 120$ expected pasta dishes.}
\end{practiceQuestion}

\begin{practiceQuestion}{9-3-q12}{mc}
\begin{questionText}
A coin is flipped $500$ times. Heads appears $260$ times. Which statement best describes this result?
\end{questionText}
\choiceA{The coin is definitely unfair.}
\choiceB{The experimental probability of $\frac{260}{500}$ is reasonably close to the theoretical probability of $\frac{1}{2}$.}
\choiceC{The theoretical probability of heads should be changed to $\frac{260}{500}$.}
\choiceD{The coin will always show heads more than tails.}
\correctAnswer{B}
\explanation{The experimental probability is $\frac{260}{500} = 0.52$, which is close to the theoretical $0.50$. Small differences between experimental and theoretical results are normal and expected.}
\end{practiceQuestion}

\begin{practiceQuestion}{9-3-q13}{mc}
\begin{questionText}
A soccer goalie blocks $27$ out of $45$ penalty kicks in practice. Based on this data, what is the experimental probability that she blocks the next penalty kick?
\end{questionText}
\choiceA{$\dfrac{2}{5}$}
\choiceB{$\dfrac{3}{5}$}
\choiceC{$\dfrac{9}{15}$}
\choiceD{$\dfrac{27}{45}$}
\correctAnswer{B}
\explanation{Experimental probability is $\frac{27}{45} = \frac{3}{5}$. This is the best estimate based on past performance data.}
\end{practiceQuestion}

\begin{practiceQuestion}{9-3-q14}{mc}
\begin{questionText}
Kai draws a card from a standard deck, records the suit, and replaces it. After $200$ trials, he records: Hearts $= 45$, Diamonds $= 53$, Clubs $= 51$, Spades $= 51$. Which suit's experimental probability is farthest from its theoretical probability?
\end{questionText}
\choiceA{Hearts}
\choiceB{Diamonds}
\choiceC{Clubs}
\choiceD{Spades}
\correctAnswer{A}
\explanation{The theoretical probability for each suit is $\frac{1}{4} = \frac{50}{200}$. Hearts appeared $45$ times, which is $5$ away from $50$. Diamonds appeared $53$ ($3$ away), Clubs $51$ ($1$ away), Spades $51$ ($1$ away). Hearts is farthest from expected.}
\end{practiceQuestion}

\begin{practiceQuestion}{9-3-q15}{mc}
\begin{questionText}
A weather station recorded rain on $54$ out of $180$ days last spring. Based on this data, how many rainy days should be expected in the next $60$ days?
\end{questionText}
\choiceA{$54$}
\choiceB{$24$}
\choiceC{$18$}
\choiceD{$30$}
\correctAnswer{C}
\explanation{Experimental probability of rain is $\frac{54}{180} = \frac{3}{10}$. In $60$ days: $\frac{3}{10} \times 60 = 18$ rainy days expected.}
\end{practiceQuestion}

% --- Short Answer Questions (6) ---

\begin{practiceQuestion}{9-3-q16}{sa}
\begin{questionText}
A bus arrives on time $168$ out of $210$ days. What is the experimental probability that the bus arrives on time? Express your answer as a simplified fraction.
\end{questionText}
\correctAnswer{$\frac{4}{5}$}
\explanation{Experimental probability is $\frac{168}{210} = \frac{4}{5}$.}
\end{practiceQuestion}

\begin{practiceQuestion}{9-3-q17}{sa}
\begin{questionText}
In a science experiment, a plant grows at least $2$ cm in $32$ out of $40$ weeks. Based on this data, what is the experimental probability (as a decimal) that the plant grows at least $2$ cm in any given week?
\end{questionText}
\correctAnswer{$0.8$}
\explanation{Experimental probability is $\frac{32}{40} = \frac{4}{5} = 0.8$.}
\end{practiceQuestion}

\begin{practiceQuestion}{9-3-q18}{sa}
\begin{questionText}
A juggler successfully completes a trick $84$ out of $120$ attempts. Based on this data, in $200$ future attempts, how many times should he expect to complete the trick successfully?
\end{questionText}
\correctAnswer{$140$}
\explanation{Experimental probability is $\frac{84}{120} = \frac{7}{10}$. In $200$ attempts: $\frac{7}{10} \times 200 = 140$ successful completions.}
\end{practiceQuestion}

\begin{practiceQuestion}{9-3-q19}{sa}
\begin{questionText}
A quality inspector tests $600$ phone chargers and finds $18$ are defective. Based on this data, how many defective chargers should be expected in a batch of $2{,}000$?
\end{questionText}
\correctAnswer{$60$}
\explanation{Experimental probability of defective is $\frac{18}{600} = \frac{3}{100}$. In $2{,}000$: $\frac{3}{100} \times 2000 = 60$ defective chargers.}
\end{practiceQuestion}

\begin{practiceQuestion}{9-3-q20}{sa}
\begin{questionText}
A student flips two coins $150$ times and gets two heads $42$ times. What is the experimental probability of getting two heads? Express your answer as a simplified fraction.
\end{questionText}
\correctAnswer{$\frac{7}{25}$}
\explanation{Experimental probability is $\frac{42}{150} = \frac{7}{25}$.}
\end{practiceQuestion}

\begin{practiceQuestion}{9-3-q21}{sa}
\begin{questionText}
A traffic engineer counts that $126$ out of $350$ cars turn left at an intersection. What is the experimental probability (as a simplified fraction) that a car turns left?
\end{questionText}
\correctAnswer{$\frac{9}{25}$}
\explanation{Experimental probability is $\frac{126}{350} = \frac{9}{25}$.}
\end{practiceQuestion}

% --- Graphical Multiple Choice Questions (3) ---

\begin{practiceQuestion}{9-3-q22}{gmc}
\begin{questionText}
The bar graph below shows the results of spinning a spinner $200$ times.

\begin{center}
\begin{tikzpicture}
  \begin{axis}[
    ybar, bar width=20pt, ymin=0, ymax=80,
    ylabel={Frequency},
    symbolic x coords={Red, Blue, Green, Yellow},
    xtick=data,
    nodes near coords,
    width=8cm, height=6cm,
    enlarge x limits=0.25
  ]
  \addplot[fill=funBlue!70] coordinates {(Red,55) (Blue,48) (Green,52) (Yellow,45)};
  \end{axis}
\end{tikzpicture}
\end{center}

Based on the data, which color has the highest experimental probability?
\end{questionText}
\choiceA{Red}
\choiceB{Blue}
\choiceC{Green}
\choiceD{Yellow}
\correctAnswer{A}
\explanation{Red appeared $55$ times out of $200$ spins, the most of any color. Its experimental probability is $\frac{55}{200} = \frac{11}{40}$, which is the highest.}
\end{practiceQuestion}

\begin{practiceQuestion}{9-3-q23}{gmc}
\begin{questionText}
The frequency table below shows the results of rolling a number cube $180$ times.

\begin{center}
\begin{tabular}{|c|c|c|c|c|c|c|}
\hline
\textbf{Outcome} & $1$ & $2$ & $3$ & $4$ & $5$ & $6$ \\
\hline
\textbf{Frequency} & $32$ & $28$ & $31$ & $30$ & $29$ & $30$ \\
\hline
\end{tabular}
\end{center}

The theoretical probability of rolling any number is $\dfrac{1}{6} \approx 0.167$. Which outcome's experimental probability is closest to the theoretical probability?
\end{questionText}
\choiceA{$1$}
\choiceB{$2$}
\choiceC{$4$}
\choiceD{$6$}
\correctAnswer{C}
\explanation{The expected frequency is $\frac{180}{6} = 30$. Outcome $4$ appeared $30$ times, exactly matching the expected value. Its experimental probability $\frac{30}{180} = \frac{1}{6}$ equals the theoretical probability.}
\end{practiceQuestion}

\begin{practiceQuestion}{9-3-q24}{gmc}
\begin{questionText}
A student tracked daily high temperatures for $100$ days and recorded whether the temperature exceeded $80^\circ$F. The results are shown below.

\begin{center}
\begin{tikzpicture}
  \fill[funRed!60] (0,0) -- (0:2.5) arc (0:230.4:2.5) -- cycle;
  \fill[funBlue!60] (0,0) -- (230.4:2.5) arc (230.4:360:2.5) -- cycle;
  \draw[thick] (0,0) circle (2.5);
  \draw[thick] (0,0) -- (0:2.5);
  \draw[thick] (0,0) -- (230.4:2.5);
  \node at (115:1.5) {\textbf{Above $80^\circ$F}};
  \node at (115:1.0) {$64$ days};
  \node at (295:1.4) {\textbf{$80^\circ$F}};
  \node at (295:0.85) {\textbf{or below}};
  \node at (295:0.3) {$36$ days};
\end{tikzpicture}
\end{center}

Based on this data, what is the experimental probability that a day's high temperature does NOT exceed $80^\circ$F?
\end{questionText}
\choiceA{$\dfrac{16}{25}$}
\choiceB{$\dfrac{9}{25}$}
\choiceC{$\dfrac{9}{16}$}
\choiceD{$\dfrac{32}{50}$}
\correctAnswer{B}
\explanation{There were $36$ days at or below $80^\circ$F out of $100$ total days. The experimental probability is $\frac{36}{100} = \frac{9}{25}$.}
\end{practiceQuestion}

% --- Graphical Short Answer Questions (3) ---

\begin{practiceQuestion}{9-3-q25}{gsa}
\begin{questionText}
The table below shows the results of a survey asking $150$ students their favorite after-school activity.

\begin{center}
\begin{tabular}{|l|c|}
\hline
\textbf{Activity} & \textbf{Number of Students} \\
\hline
Sports & $54$ \\
\hline
Video Games & $42$ \\
\hline
Reading & $24$ \\
\hline
Art & $30$ \\
\hline
\end{tabular}
\end{center}

Based on this data, what is the experimental probability that a randomly selected student prefers reading? Express your answer as a simplified fraction.
\end{questionText}
\correctAnswer{$\frac{4}{25}$}
\explanation{There are $24$ students who prefer reading out of $150$ total. The experimental probability is $\frac{24}{150} = \frac{4}{25}$.}
\end{practiceQuestion}

\begin{practiceQuestion}{9-3-q26}{gsa}
\begin{questionText}
The double bar graph below compares experimental results (dark bars) with the expected theoretical results (light bars) for spinning a spinner $240$ times.

\begin{center}
\begin{tikzpicture}
  \begin{axis}[
    ybar, bar width=12pt, ymin=0, ymax=100,
    ylabel={Frequency},
    symbolic x coords={Red, Blue, Green, Yellow},
    xtick=data,
    width=9cm, height=6cm,
    enlarge x limits=0.2,
    legend style={at={(0.5,1.05)}, anchor=south, legend columns=2}
  ]
  \addplot[fill=funBlue!80] coordinates {(Red,68) (Blue,55) (Green,62) (Yellow,55)};
  \addplot[fill=funBlue!25] coordinates {(Red,60) (Blue,60) (Green,60) (Yellow,60)};
  \legend{Experimental, Theoretical}
  \end{axis}
\end{tikzpicture}
\end{center}

Which color's experimental probability differs the most from the theoretical probability? By how much do they differ? Express the difference as a simplified fraction.
\end{questionText}
\correctAnswer{Red; $\frac{1}{30}$}
\explanation{The theoretical frequency for each color is $60$ out of $240$. Red appeared $68$ times — a difference of $8$. The difference in probability is $\frac{8}{240} = \frac{1}{30}$, which is the largest deviation among the four colors.}
\end{practiceQuestion}

\begin{practiceQuestion}{9-3-q27}{gsa}
\begin{questionText}
An archer shoots $80$ arrows at a target and records hits and misses over $4$ practice sessions.

\begin{center}
\begin{tabular}{|l|c|c|}
\hline
\textbf{Session} & \textbf{Hits} & \textbf{Misses} \\
\hline
Session 1 & $14$ & $6$ \\
\hline
Session 2 & $16$ & $4$ \\
\hline
Session 3 & $15$ & $5$ \\
\hline
Session 4 & $17$ & $3$ \\
\hline
\end{tabular}
\end{center}

What is the overall experimental probability that the archer hits the target? Express your answer as a simplified fraction.
\end{questionText}
\correctAnswer{$\frac{31}{40}$}
\explanation{Total hits: $14 + 16 + 15 + 17 = 62$. Total arrows: $80$. The experimental probability is $\frac{62}{80} = \frac{31}{40}$.}
\end{practiceQuestion}
